\addcontentsline{toc}{section}{\MakeUppercase{Santrauka}}
\noindent\begin{minipage}{\textwidth}
	\textbf{\MakeUppercase{Santrauka}}
    \setcounter{section}{0}
    % \addtocontents{toc}{\protect\setcounter{tocdepth}{0}}
	\subsection*{\k{I}vadas}\label{chap:introlt}
    % \addtocontents{toc}{\protect\setcounter{tocdepth}{2}}
\end{minipage}
\subsection*{Darbo aktualumas}
Garso klasifikavimas atlieka svarbų vaidmenį aptinkant anomalijas, nes suteikia vertingos informacijos įvairiais taikymo atvejais~\cite{Kilickaya2025, LoScudo2023, Barusco2025}. Šios sistemos yra naudingos pramoninėse ir miesto vietovėse~\cite{Purohit2019, Piczak2015, Salamon2014}, taip pat sudėtingose gamtinėse vietovėse, pvz., miškuose~\cite{Bandara2023}. Garso duomenys gali būti \dots

\dots
\subsection*{Problemų formulavimas}
\begin{enumerate}[leftmargin=*]
	\item \dots
	\item \dots
\end{enumerate}
\subsection*{Tyrimo tikslas}
Šio tyrimo tikslas – pagerinti garso pagrindu veikiančių anomalijų aptikimo sistemų tikslumą, patikimumą ir \dots

\subsection*{Tyrimo uždaviniai}
\begin{enumerate}[leftmargin=*]
	\item \dots
	\item \dots
\end{enumerate}

\subsection*{Mokslinės naujovės}
Šiame tyrime pristatomos šios mokslinės naujovės, kurios išplečia \dots

\begin{enumerate}[leftmargin=*]
	\item \dots
	
	\item \dots
\end{enumerate}

\subsection*{Praktinė reikšmė}
\begin{enumerate}[leftmargin=*]
	\item \dots
	\item \dots
\end{enumerate}

\subsection*{Disertacijos teiginiai}
\begin{enumerate}
	\item \dots
	\item \dots
\end{enumerate}

\subsection*{Moksliniai aprobavimai}
Šio disertacinio tyrimo metodologiją ir \dots

\subsection*{Disertacijos struktūra}

Disertaciją sudaro penki skyriai. \hyperref[chap:introlt]{„Įvadas“} skyriuje pristatoma tyrimo problema ir jos motyvacija, apibrėžiama šiame darbe vartojama akustinių anomalijų sąvoka ir pagrindžiamas pramoninių, urbanistinių bei miško garso aplinkų pasirinkimas kaip papildomų vertinimo sričių. \hyperref[chap:twolt]{„Literatūros apžvalgos apie giluminio neuroninio tinklo metodą akustinių anomalijų įvykiams aptikti“} skyriuje pateikiama \dots
